% from ega0-19-fr.tex
\oldpage[0\textsubscript{IV-1}]{84}

\begin{proposition}[19.3.10]
\label{0.19.3.10-fr}
Soient $A$ un anneau topologique, $B$ une $A$-algèbre formellement lisse.
Soient $C$ une $A$-algèbre topologique, $\mathfrak J$ un idéal de $C$,
vérifiant les conditions suivantes~:
\begin{enumerate}
  \item[$1^\circ$] $C$ est métrisable et complet.
  \item[$2^\circ$] $\mathfrak J$ est fermé et la suite
  $(\mathfrak J^n)$ tend vers $0$.
\end{enumerate}
Alors, tout $A$-homomorphisme continu $u:B\to C/\mathfrak J$ se factorise en
$B\xrightarrow{v}C\to C/\mathfrak J$, où $v$ est un
$A$-homomorphisme continu.
\end{proposition}

\begin{proof}
Soit $(\mathfrak Q_n)$ une suite décroissante d'idéaux de $C$ formant un
système fondamental de voisinages de $0$. Pour tout $n$, considérons la
$A$-algèbre discrète $C/\mathfrak Q_n$ et l'idéal
$(\mathfrak J+\mathfrak Q_n)/\mathfrak Q_n$ de cette algèbre~; comme il
existe $k$ tel que $\mathfrak J^k\subset\mathfrak Q_n$,
$(\mathfrak J+\mathfrak Q_n)/\mathfrak Q_n$ est nilpotent. Pour tout
$n$, soit $u_n$ l'homomorphisme continu
\[
  B\xrightarrow{u}C/\mathfrak J\to C/(\mathfrak J+\mathfrak Q_n)
  =(C/\mathfrak Q_n)/((\mathfrak J+\mathfrak Q_n)/\mathfrak Q_n)~;
\]
par hypothèse $u_n$ se factorise en
\[
  B\xrightarrow{v_n}C/\mathfrak Q_n
  \xrightarrow{\varphi_n}C/(\mathfrak J+\mathfrak Q_n),
\]
où $v_n$ est un $A$-homomorphisme continu. Montrons qu'on peut choisir de
proche en proche les $v_n$ de sorte que $v_n$ se factorise en
\[
  B\xrightarrow{v_{n+1}}C/\mathfrak Q_{n+1}\to C/\mathfrak Q_n
\]
pour tout $n$ (autrement dit, que les $v_n$ forment un \emph{système
projectif} d'homomorphismes). Cela résultera du lemme suivant.
\end{proof}

\begin{lemma}[19.3.10.1]
\label{0.19.3.10.1-fr}
Soient $B$ une $A$-algèbre formellement lisse, $E$, $E'$ deux $A$-algèbres
topologiques discrètes, $\mathfrak R$ (resp. $\mathfrak R'$) un idéal
nilpotent de $E$ (resp. $E'$), $f:E\to E'$ un $A$-homomorphisme surjectif tel
que $f(\mathfrak R)\subset\mathfrak R'$, $f':E/\mathfrak R\to
E'/\mathfrak R'$ l'homomorphisme déduit de $f$ par passage aux quotients.
Soient $u:B\to E/\mathfrak R$ un $A$-homomorphisme continu, $u'$
l'homomorphisme composé
$B\xrightarrow{u}E/\mathfrak R\xrightarrow{f'}E'/\mathfrak R'$, et soit
$v':B\to E'$ un $A$-homomorphisme continu tel que $u'$ se factorise en
$B\xrightarrow{v'}E'\to E'/\mathfrak R'$. Alors il existe un
$A$-homomorphisme continu $v:B\to E$ tel que $v'$ se factorise en
$B\xrightarrow{v}E\xrightarrow{f}E'$.
\end{lemma}

Dans le diagramme commutatif
\[
\xymatrix{
  E \ar[r]^f \ar[d] & E' \ar[d] \\
  E/\mathfrak R \ar[r]_{f'} & E'/\mathfrak R'
}
\]
tous les homomorphismes sont surjectifs~; si $\mathfrak L$ est le noyau de
$f$, $E'$ s'identifie à $E/\mathfrak L$ et $\mathfrak R'$ à
$(\mathfrak R+\mathfrak L)/\mathfrak L$. Utilisons maintenant le lemme
élémentaire suivant.

\begin{lemma}[19.3.10.2]
\label{0.19.3.10.2-fr}
Soient $F$ un anneau (non nécessairement commutatif), $\mathfrak a$,
$\mathfrak b$ deux idéaux bilatères de $F$~; alors le produit fibré
\[
  (F/\mathfrak a)\times_{F/(\mathfrak a+\mathfrak b)}(F/\mathfrak b)
\]
s'identifie canoniquement à $F/(\mathfrak a\cap\mathfrak b)$.
\end{lemma}

\begin{proof}
Ce n'est autre qu'une application particulière de (18.1.7), où, dans le
diagramme (18.1.7.1), on remplace $G$ par
$F/(\mathfrak a\cap\mathfrak b)$, $\mathfrak J'$ par
$\mathfrak b/(\mathfrak a\cap\mathfrak b)$, $F$ par
$F/\mathfrak b$, $\mathfrak R$ par
$(\mathfrak a+\mathfrak b)/\mathfrak b$, $B$ par
$F/(\mathfrak a+\mathfrak b)$ et $\theta$ par le $F$-homomorphisme
bijectif canonique
$(\mathfrak a+\mathfrak b)/\mathfrak b\to
\mathfrak a/(\mathfrak a\cap\mathfrak b)$.
\end{proof}

